\part Software {{\bf Allen's Law:} Everything is more complicated
than it appears.}

\chap Software Specification

This chapter outlines the user interface and the functions to be performed
by the software model of a back-propagation \nn.
In the following specification, a {\sl message\/} will convey some
information on program operation to the user, a {\sl warning\/} is a
non-catastrophic signalling to the user, while an {\sl error\/} will
cause a termination of the program after an error message is printed.

The specification as it is presented here has been fully implemented in
the software package, called fishNET.
The design used to implement this specification is described in
Chapter~4: Construction, and its behaviour is examined in Chapter~5:
Performance.

\sec The User Interface

Operation of the program requires a number of parameters and fishNET
offers a number of different ways of saving and viewing the
data about network functions and the results of learning.
The program accepts command line inputs of names of files which
contain network data, or can accept input from the keyboard if a
data file is absent.
Both methods of input (file and keyboard) will be described below.

\ssec Selecting the Parameters -- Configuring the Network

The user is able to select the size and shape of the network, as
well as the parameters for use during learning.
Entry is via an input file specified on the command line, or by
answering questions about the individual data items if no command
line file name is given.

The information needed to run the program is:
\smallskip
\item{$\bullet$} A comment or description of the problem for saving
with the network.
\item{$\bullet$} The number of layers in the network, and the number
of \neus\ in those layers.
\item{$\bullet$} The names of the files containing teaching input
data and expected data.
\item{$\bullet$} The name of the file containing data for use in
execution and the name of the output file.
\item{$\bullet$} The name of the file where the network is to be
saved.
\item{$\bullet$} The learning parameters~$\alpha$ and~$\varepsilon$.
\item{$\bullet$} The nominal output width of the output layer of
\neus\ for output formatting.
\item{$\bullet$} The maximum number of input/output pair learning
sweeps required, if a limit is desired.
\smallskip
If no command line file name is specified, the program will query
the user via the keyboard.
This is the default operation.

\sssec Keyboard Input Format

A question is asked via the screen for every parameter listed above,
and the answers are accepted through the keyboard.
The user enters a return after typing the parameter being
considered.

\sssec Configuration File Format

The format of a configuration file is described below.
The symbols~$\langle\rangle$ are used to delimit user required data.
All information in double quotes ``'' are tokens and must be
inserted {\sl exactly} as they appear otherwise an error will occur.
(The error message is `appropriate token not found'.)
All tokens and data must also appear in the order shown below.
The exception is any piece of data which is surrounded
by~$\{\}$ brackets.
These data items are optional.
If they are used, however, please note the quoted and bracketed
tokens and data items necessary.
$$\vbox{
\halign{#\hfil&&\ \hfil#\hfil\cr
	\noalign{\noindent$\{\langle \ldots$ user comment $\ldots \rangle\}$}
	\noalign{\smallskip}
	``[start]''\cr
	\noalign{\smallskip}
	``layers''&$\langle n\rangle$\cr
	``neurons''&$\langle a_0$&$a_1$&$a_2$&$\ldots$&$a_{n-1}\rangle$\cr
	\noalign{\smallskip}
	``output width''&$\langle m\rangle$\cr
	\noalign{\smallskip}
	``alpha''&$\langle\alpha\rangle$\cr
	``epsilon''&$\langle\varepsilon\rangle$\cr
}}
$$
$$\vbox{
\halign{#\hfil&\ #\hfil\cr
	\noalign{$\quad$}
	``sample in''&$\{$path$\}\langle$teaching input file name$\rangle$\cr
	``sample out''&$\{$path$\}\langle$expected output file name$\rangle$\cr
	\noalign{\smallskip}
	``execute in''&$\{$path$\}\langle$test data file name$\rangle$\cr
	``execute out''&$\{$path$\}\langle$test output file name$\rangle$\cr
	\noalign{\smallskip}
	``network save''&$\{$path$\}\langle$network save file name$\rangle$\cr
	\noalign{\smallskip}
	$\{$``max sweeps''&$\langle q\rangle\}$\cr
}}
$$

Simply place the desired value next to the appropriate token.
Tokens must be in order.
For the ``neurons'' token, the order of the numbers~$a_0\to a_{n-1}$ is 
taken as the number of \neus\ in the layers~$0\to n-1$.
Hence the example
$$\vbox{
\halign{#\hfil&&\ \hfil#\hfil\cr
	\noalign{$\quad\vdots$}
	layers&$3$\cr
	\neus &$13$&$20$&$5$\cr
}}
$$
means that there are 13 neurons on the input (layer~$0$),~$20$ in
the intermediate layer (layer~$1$), and~$5$ in the output layer
(layer~$2$).

The tokens ``alpha'' and ``epsilon'' refer to the learning
parameters $\alpha$ and $\varepsilon$.

The file name following the ``sample in'' token is the file used as
training input for the network, and the file name following the
``sample out'' token is the expected output file, used for
comparison and calculating the error.

The ``execute in'' and ``execute out'' tokens precede the file names
for use in operating the network once it is taught and the
final destination of the output.

The token ``network save'' comes before the name of the file where
the network is to finally be saved once it has been taught, or the
program terminated.

Finally, ``max sweeps'' is the maximum number of training sweeps
that will occur.
Once this number is reached, (and it will not be reached if the
network is completely taught,) the network is saved and
execution begins.

Formats of data files are shown below.

Command line options used to specify a configuration file instead of
keyboard entry are found in section~3.1.5.\ Command line options.

\ssec Using a Pre-Made Network

The user is able to load a network file which contains all the data
necessary (especially weights) for operation of
the network, or continued teaching, or both.
The code is designed so that if it is halted (a Ctrl-c in MSDOS, or a
BREAK or DEL in Unix) the network's present status and elapsed~$t$ are
stored into a network file.
The network is also saved when it has been fully taught, or the
maximum time has elapsed (depending upon flags set, so see~3.1.5.\
Command line options).

\sssec Network File Format

The format of a network file is described below.
The symbol format is the same as the previously described
configuration file format (3.1.1.2).
$$\vbox{
\halign{#\hfil&&\ \hfil#\hfil\cr
	\noalign{\noindent$\{\langle \ldots$ user comment $\ldots \rangle\}$}
	\noalign{\smallskip}
	``[start]''\cr
	\noalign{\smallskip}
	``layers''&$\langle n\rangle$\cr
	``neurons''&$\langle a_0$&$a_1$&$a_2$&$\ldots$&$a_{n-1}\rangle$\cr
	``weights''&$\langle b_0$&$b_1$&$b_2$&$\ldots$&$b_{m-1}\rangle$\cr
	\noalign{\smallskip}
	``output width''&$\langle m\rangle$\cr
	\noalign{\smallskip}
	``alpha''&$\langle\alpha\rangle$\cr
	``epsilon''&$\langle\varepsilon\rangle$\cr
}}
$$
$$\vbox{
\halign{#\hfil&\ #\hfill\cr
	\noalign{$\quad$}
	``sample in''&$\{$path$\}\langle$teaching input file name$\rangle$\cr
	``sample out''&$\{$path$\}\langle$expected output file name$\rangle$\cr
	\noalign{\smallskip}
	``execute in''&$\{$path$\}\langle$test data file name$\rangle$\cr
	``execute out''&$\{$path$\}\langle$test output file name$\rangle$\cr
	\noalign{\smallskip}
	``start time''&$\langle q\rangle$\cr
	``learn time''&$\langle r\rangle$\cr
}}
$$

Place the desired value next to the appropriate token.
Tokens must be in order.
For the ``\neus'' token, the order of the numbers $a_0\to a_{n-1}$ is
taken as the same as the number of \neus\ in the layers~$0\to n-1$.
Data following the ``weights'' token behaves identically: if there
are~$a_0$ \neus\ in layer~$0$, and~$a_1$ in layer~$1$, then there
will be $a_0 \times a_1$ weights between layer~$0$ and layer~$1$.
The first $a_1$ weights are the strengths of connections between
the first \neu\ in layer~$0$ and the \neus\ in the layer above.
The second $a_1$ weights apply to the second \neu's connections to
the layer above, etc., until the $a_0$th $a_1$ weights apply to the
last \neu\ in layer~$0$.
The pattern is repeated for all layers (note that there are no weights
associated with the output $n-1$th layer).

For example, consider the following fragment of a network file:
$$\vbox{
\halign{#\hfil&&\ \hfil#\hfil\cr
	\noalign{$\vdots$}
	``layers''&$3$\cr
	``neurons''&$1$&$2$&$3$\cr
	``weights''&$w_{1,1}$&$w_{1,2}$&$w_{2,1}$&$w_{2,2}$&
		$w_{2,3}$&$w_{2,4}$&$w_{2,5}$&$w_{2,6}$\cr
	\noalign{$\vdots$}
}}
$$
Here there are~$3$ layers; the input layer contains~$1$ \neu, the
middle layer contains~$2$ \neus, and the top layer~$3$ \neus.
The weights $w_{1,1}$ and $w_{1,2}$ connect the single input \neu\ to the
first and second of the~$2$ second layer \neus\ respectively.
The weights $w_{2,1}\to w_{2,3}$ connect the first \neu\ in the
second layer with the \neus~$1\to 3$ in the top layer, etc.

The token ``start time'' marks the value of $t$ that the network
will start learning from, if the user desires continued teaching of
the network.
In other words, it is a measure of how long the network had been
taught for if it was generated by a Ctrl-c or BREAK, or saved when
it had finished learning or reached the ``max sweeps'' value of
$t$.

The token ``learn time'' is equivalent to the ``max sweeps'' token
and indicates to the software what value of $t$ to stop 
training the network at if the user desires continued teaching of
the network.

Command line options used to specify a network file and what is to
be done with it (that is continued teaching or execution) are found
in section 3.1.5.~Command line options.

\ssec Teaching Data

The specifications for configuration and network files, and
keyboard input contain references to ``sample in'' and
``sample out'' files.
These refer to the teaching data, used by the program to teach the
network.

The ``sample in'' file contains the data patterns or values desired
for learning, and the ``sample out'' file contains the expected
output for the sample input.
Input/output pairs are matched; the first data items in the input
file will produce the first data items in the output file after
learning.
The second set of input data will match the second set
of output data after learning, and so on.
Each input output pair is called a {\sl case}.

\sssec Teaching Data File Format

The start of an input case is signalled by a ``[start]'' token.
There may be a user comment before the first token.
Input data is matched to the input \neus\ one at a time, and in the
order that they appear in the file.
If there are more data items than there are \neus, the excess ones
are thrown away up until the next ``[start]'' token occurs, and a
warning is issued.
The next case is then parsed.
If for any case there are more input \neus\ than there is data, an
error occurs.

The format for the file is as shown below:
$$\vbox{
\halign{\hfil#\hfil&&\ \hfil#\hfil\cr
	\noalign{\noindent$\{\langle\ldots$ user comment $\ldots\rangle\}$}
	\noalign{\noindent``[start]''}
	$\alpha_0$&$\beta_0$&$\gamma_0$&$\delta_0$&$\ldots$\cr
	\noalign{\noindent``[start]''}
	$\alpha_1$&$\beta_1$&$\gamma_1$&$\delta_1$&$\ldots$\cr
	\noalign{\noindent``[start]''}
	$\alpha_2$&$\beta_2$&$\gamma_2$&$\delta_2$&$\ldots$\cr
	\noalign{\quad$\vdots$}
}}
$$

\sssec Expected Data File Format

Expected data files have the same format as teaching data files.
The values in the expected data file are matched to output \neus\ in
the same way as input data is matched to the input layer \neus.
If for any case there are more data items than there are \neus, the excess ones
are thrown away up until the next ``[start]'' token occurs, and a
warning is issued.
If for any case there are more output \neus\ than there is output
data, an error occurs.
If there is not an identical number of input and output cases, an
error occurs.

\ssec Test Data

Configuration and network files, and keyboard input, contain the
tokens ``execute in'' and ``execute out''.
The file names after these tokens refers to the files used as input
and output for the network after it is trained.
The ``execute in'' file contains data for the inputs of the bottom
(layer~$0$) \neus, and the ``execute out'' file is used to store
the output from the top (layer~$n-1$) \neus\ for each case presented to
the bottom layer.
Any previous contents of the ``execute out'' file are overwritten by
the new input cases' outputs.

\sssec Test Data File Format

Format is identical to 3.1.3.1.~Teaching data file format.
Once again, if there is too much data the excess for each input
cases is ignored (and a warning is issued), and if there is too
little data an error occurs.

\sssec Output Data File Format

The ``execute out'' file output format is similar to 3.1.3.2.~Expected
data file format.
The only difference is that the length of the stored output line is
modulated by the parameter following the token ``output width'',
which is entered from the keyboard, or network or configuration
files.

\ssec Command Line Options

To reduce keyboard input and increase the facility for repeated
(automated) testing of program behaviour, all directives and flags
used by the program can be set from the command line.
Each flag or directive is selected by typing
$$
{\rm fishNET\ } \{-flag_0\}\{-\}\{flag_1\}\ldots\quad.
$$
For example, to specify the flags $-x$ and $-v$ (the meaning of
these two flags will be explained in detail below,) you could type
${\rm fishNET\ } -x {\rm \thinspace} -v$ or ${\rm fishNET\ } -xv$.
The exception to this method are the $-c$ and $-n$ flags, which can
only appear at the end of a ``$-$'' list or by themselves.
This is because they refer to file names.
For example,
$$
\hbox{fishNET\ $-n\langle{\rm netfile}\rangle$\ $-q$.}
$$

A full explanation of each flag and directive appears below
categorised by function.

\sssec Screen Input/Output Control

\ssssec $-v$ (verbose) flag

The $-v$ flag causes the program to produce verbose runtime
information about the present status of the network.
It displays the expected output from the network alongside the
actual output, a running measure of errors, and the elapsed time,
$t$.
This process slows down execution speed by the amount of time taken
for screen output.

This flag cannot be used with the $-q$ flag.

\ssssec $-q$ (quiet) flag

The $-q$ flag causes suppression of all runtime messages.
This flag is good for batch mode processing, and of course gives the
fastest execution speed by virtue of minimal screen output.

This flag cannot be used with the $-v$ flag.

\ssssec Default

If neither the $-q$ or $-v$ flags are set, the default mode is used.
This mode causes fishNET to print rudimentary information 
about loading status, and
displays the elapsed time and total error while the program is
running.

\sssec Loading Configuration and Network Files from Disk

\ssssec $-c$ (load configuration file) directive

The $-c\{{\rm path}\}\langle{\rm file name}\rangle$ directive lets
the user load a configuration file of parameters into the simulator.
The format of the file is described in an earlier section entitled
3.1.1.2.~Configuration file format.

This directive can't be used with the $-n$ directive.

\ssssec $-n$ (load network file) directive

The $-n\{{\rm path}\}\langle{\rm file name}\rangle$ directive enables
the user to load a pre-taught or manually created network file into the
simulator.
The format of the file is described in an earlier section, 
3.1.2.1.~Network file format.

This directive can't be used with the $-c$ directive.

\ssssec Default

There is no default file loaded.
If neither $-c$ or $-n$ is specified, the user is queried via the
screen and keyboard for the appropriate parameter values.

\sssec Using the Network File

\ssssec $-e$ (execute) flag

The $-e$ flag is used after the $-n\langle{\rm name}\rangle$
directive to tell the simulator to execute the network with the data
files whose names are included in the network file.
The network will {\sl not} continue to be trained.

The $-e$ flag can't be used with the $-t$ flag.

\ssssec $-t$ (teach) flag

The $-t$ flag is used after the $-n\langle{\rm name}\rangle$
directive to tell the simulator to continue teaching the network
with the sample input and output files whose names are included in
the network file.
When the network is ``taught'', execution with the data file (the
name of which is also in the network file,) will occur.

The $-t$ flag can't be used with the $-e$ flag.

\ssssec Default

If neither the $-e$ or $-t$ flags are used with the $-n$ directive,
an error occurs.
In other words, you must tell the program what you want to do with a
network file!

\sssec Saving the Network File

\ssssec $-s$ (store ``taught'' network) flag

The $-s$ flag tells the program to store the network after it is
taught.
If this flag is set, once the errors drop below the threshold, the
network is automatically stored in the file named through the
keyboard, network, or configuration files.
The format of the saved file is specified in an earlier section,
3.1.2.1.~Network file format.

This flag can't be used with the $-e$ (execute) flag, or the $-d$ or
$-x$ flags (see below).

\ssssec $-d$ (don't store network) flag

The $-d$ flag indicates that once the network is taught, it will not
be stored.
If this flag is set, execution begins immediately after the network
is taught, and although the network is not saved to disk, the output
from the execution is.

This flag can't be used with the $-s$ or $-x$ save flags. If it is
used with the $-e$ flag it is ignored.

\ssssec $-x$ (store learning information and time) flag

The $-x$ flag is intended primarily for experimental use of the
program.
When this flag is set, the program will teach or continue to teach
the network but when it is ``taught'' only the necessary parameters
are saved to disk.
These parameters are the weight modification factors~$\alpha$
and~$\varepsilon$, the number of layers and \neus\ in those layers,
I/O files and time required to teach the network.
This saves time and disk space for users of fishNET who are
experimenting with the parameters' and data's effects on
learning speed.

This flag can't be used with the $-s$ or $-d$ save flags, or the
$-e$ execute flag.

\ssssec Default

If none of the save flags is specified, the user is queried via the
screen and keboard. The choices available are akin to the $-s$ and
$-d$ flags (there is no equivalent $-x$ option).

\sssec Calculation of $\Delta w(t)$

The package fishNET offers two ways to calculate $\partial E \over
\partial w$ and $\Delta w(t)$, both enumerated in algorithmic form
in the previous chapter.

\ssssec $-1$ (each) mode

When the $-1$ flag is set, the program will calculate and apply
$\Delta w(t)$ for each I/O case.
This mode is only included for experimental purposes (for a
description of the algorithm used, see the second algorithm in the
previous chapter).

The $-1$ flag can't be used with the $-a$ flag.

\ssssec $-a$ (all) mode

If the $-a$ flag is specified, the calculation of $\Delta w(t)$ is
once per sweep (presentation) of all I/O cases.
$\partial E \over \partial w$ values are calculated for each I/O
case and added together.
$\Delta w(t)$ is calculated at the end of each sweep from the
accumulated $\partial E \over \partial w$ value.
This is the most commonly used mode.

The $-a$ flag can't be used with the $-1$ flag.

\ssssec Default

If neither the $-1$ or $-a$ flags are specified, the $-a$ mode is
assumed.
Unless $-q$ (quiet) has been set, a message is issued to this
effect.

\sssec Printing the Help Message

\ssssec $-?$, $-h$ (help) flags

If the $-?$ or $-h$ flags are specified, the program doesn't run and
a help message is displayed.
The help message contains information on every flag and is located
in the file {\tt help.hlp} under all implementations of the program.

\sssec Unrecognised Options

If any other option is specified, the program displays the message 
$$
\hbox{``try typing `fishNET $-?$' for help''.}
$$

\sssec Default (no options set)

If no options are set, the program asks for all information about
file names and parameters, the same way it does if neither $-c$ or
$-n$ is specified.

This is the recommended mode for ``beginners''.

\sec Operations Performed

FishNET implements the back-propagation model of a
\nn, as described in the previous chapter ``The Back-Propagation
Model''.
The shape of the network is flexible and specifiable by the user,
and as many parameters as possible are easily modifiable.

The most important parameters are the network shape and size,
and the initial values of~$\alpha$ and~$\varepsilon$.
These values are user specifiable, and used during learning.

Input and sample output files are user specifiable and used for
input/output pairs.

Final network operation output file is user specifiable.

Efficiency of operation is important, however the operations required
are undertaken in such a way that flexibility of network size and
portability of code are maintained.

\sec Network Size

Any user specified network can have up to $8$ layers.
It can have any number of \neus\ in these layers, depending solely upon
available memory of the particular hardware implementation.

Allocation of space for the network is completely dynamic.

\sec File Input/Output Format

All files used by and created by the program are ASCII files.
As a direct result, data files can be generated by hand with an
editor, as can network and configuration files.
